\chapter{Évaluer les modèles d'IA}
\label{ch:jour5}

\begin{aitip}
Ce chapitre est en cours de rédaction. Le cours est esquissé ci-dessous ; le texte complet sera publié dans la prochaine itération de l'atelier. Le syllabus et le notebook de TP de cette journée sont déjà disponibles.
\end{aitip}

\section*{Sections prévues}
\begin{itemize}
  \item Statistiques de base, décomposition biais--variance, signification du sur-apprentissage.
  \item Stratégies de validation, calibration, quantification de l'incertitude.
  \item Exemples adversariaux et données hors distribution --- pourquoi un modèle parfait en laboratoire échoue en production.
  \item Le cadre de décision : quand déployer un modèle, quand collecter plus de données, quand reformuler le problème.
\end{itemize}

\section*{Travail pratique}
Le notebook du Jour 5 (\texttt{Day5\_Evaluation\_Generalization.ipynb}) compare plusieurs architectures sous des décalages de distribution contrôlés, construisant l'intuition de quels échecs se généralisent et lesquels ne se généralisent pas.
